\begin{abstractEN}

In recent years, non-geostationary satellite orbit (NGSO) systems, especially low Earth orbit (LEO) satellite constellations, have become an important enabler for future broadband satellite communications and non-terrestrial networks. However, when LEO downlink transmissions operate in the same or adjacent frequency bands as existing geostationary satellite orbit (GSO) systems, they may cause co-channel interference to GSO ground stations. To protect GSO systems, NGSO systems must comply with equivalent power flux-density (EPFD) limits. When an LEO satellite passes near a GSO ground station and forms a near-inline geometry with respect to the GSO direction, some high-risk beams may need to be shut down to satisfy the EPFD limit. Although beam shutdown can effectively reduce interference, it also interrupts the service of users originally served by the closed beams and degrades LEO user service continuity and satisfaction.

To address this issue, this thesis proposes Safe-Beam Assisted Power-Release Relay (SAPR-R), a staged relay-based service recovery method for multi-beam LEO satellite systems under NGSO--GSO coexistence. SAPR-R exploits the predictable motion of LEO satellites and beam footprints to precompute the critical satellite, helper satellites, closed beams, safe beams, and beam overlap mappings between critical and helper satellites. When the critical satellite shuts down high-risk beams to satisfy the EPFD limit, SAPR-R first lets safe beams on the critical satellite take over selected helper users to release available power on the helper satellite; then allocates the released power to helper relay beams to improve their capability to serve closed-beam users; finally, it reassociates affected users to feasible helper relay beams through beam-pair local relay to complete service recovery.

An STK--MATLAB simulation framework is developed to evaluate the EPFD compliance and user service performance of SAPR-R during the critical satellite pass near the GSO ground station. The results show that SAPR-R keeps the aggregate EPFD below the regulatory limit and achieves higher average user satisfaction than baseline methods. Moreover, SAPR-R increases the number of closed-beam users successfully served by helper relay beams and reduces the proportion of low-satisfaction users. Overall, SAPR-R improves LEO user service continuity while maintaining EPFD compliance, without solving a high-complexity global optimization problem in every time slot, thereby providing a practical and feasible relay-based service recovery approach.

$\bf{Keywords}$ : NGSO--GSO coexistence, EPFD, LEO satellites, multi-beam satellite systems, service recovery, relay, non-terrestrial networks, SAPR-R.

\end{abstractEN}
